\documentclass[12pt]{article}

\usepackage[a4paper, total={6in, 8in}]{geometry}
\usepackage{textcomp}

\begin{document}
\setlength{\parindent}{0pt}

\def \t {\textrightarrow}
\def \v {\vspace{1em}}
\def \bi {\begin{itemize}}
\def \ei {\end{itemize}}
\def \s[#1] {\section*{#1}}
\def \ss[#1] {\subsection*{#1}}
\def \sss[#1] {\subsubsection*{#1}}

\s[Dopoguerra in Italia]
Forse bisogna recuperare lezione precedente? \\
Nel 46 Vittorio Emanuele III abdica perche cerca di mantenere l'istituzione monarchica \\
Togliatti, segretario del partito comunsita \t dice che bisogna aspettare, uscire da crisi e scegliere assetto istituzionale del paese \\
Nenni, De Gasperi, Togliatti (chiedere a kathirn questa parte, tra pov politico e 2 giugno '46) \\
Elezioni, vince la repubblica e viene eletta l'assembla costituente \t c'è pero chi dice che ci sono stati del brogli elettorali (al sud vinche la monarchia) \\
Vince la repubblica, il re va in esilio \\
Con l'assemblea costituente si afferma democrazia cristiana (maggiore successo, 207 seggi), poi socialisti (115) e comunisti (104) \\
Per l'assemblea costituente si era presentato anche "Il fronte dell'uomo qualunque" \t che dura poco, ma esprime una mentalita qualunquista \\
Il 25 giugno del 46 l'assemblea si riunisce e nomina presidente dell'assemblea socialista Saragat \t mentre il capo provvisorio è de Nicola (che poi firmera la costituzione) \\
All'assemblea partecipano tanti antifascisti, Pertini (poi presidente), Benedetto Croce etc. \\
Prima c'era statuto albertino \\
La costituente lavora molto, ed elabora costituizione (approvata nel dicembre del 47 ed entra in vigore 1 gennaio del 48)

\v

Lo scontro + duro si avra riguardo i rapporti con la chiesa \\
Sia partiti laici sia socialisti sono contrari all'introduzione dei patti laternansi, ma non i comunisti \t Togliatti voleva mantenere contatti con democrazia cristiana \\
L'articolo 7 viene approvato con il voto contrario dei socialisti e dei partiti laici \t articolo 7 viene anche inserito con la dicitura che "ogni modiifica di questo articolo puo essere fatta senza rivisitazione costit." \t cosi nell 84 Craxi firma nuovo accordo con chiesa, che sostituisci Patti Laternansi \t che fa l'italia stato laico (mentre articolo 7 faceva italia stato confessionale) \\
La pace pero non è stata ancora firmata \t quindi governo de gapseri ha problema del trattato di pace, che si fa nella conferenza di Londra \\
Italia non ha posizione forte nelle trattative \t viene ritenuta comunque responsabile nonostante la resistenza e il cambio politico nel 43 \\
De Gasperi stava sempre + costruendo un rapporto di amicizia con USA, mentre Togliatti privilegiava rapporti con URSS (ma stava lavorando per posizione equidistante dalle potenze dell italita)

\v

Pace firmata nel febbraio del 47 \t e le condizioni sono dure:
\bi
  \item perde parte venezia giulia
  \item trieste diventa citta internazionale
  \item perde tutte le colonie africane
  \item deve pagare riparazioni di guerra ai paesi aggrediti
\ei
De Gasperi va anche negli USA, e torna con un prestito di 100 milioni di dollari (pochi, atto simbolico) \t torna pero con un amicizia con USA \\
Si capisce che la d.c. è l'unica forza che puo arginare comunismo \t De Gasperi si allaccia a Truman, che fa politica anticomunista \\
De Gasperi è anche filostatunitense (ed è presidente del consiglio, quindi decide lui direzione dell'italia) \\
Collaborazione tra d.c. e sinistra è impossibile \t si allega anche con Truman, che dice che combattera il comunismo \t quindi rapporti con Togliatti si perdono

\v

Nel 47 fa nuovo governo senza sinistre \t chiamati "Anni del centrismo" \t anni dominati dalla democrazia cristiana, fino al 1962 \\
Nel 1949 Italia entra nella NATO = alleanza militare che nasce nel 49, di cui fanno parte paesi europei e USA e Canada \\
Clima in italia non è disteso \t guerra fredda, italia schierata \t cosi nell aprile del 48 si tengono elezioni con nuova costituzione, dove d.c. ottiene maggioranza assoluta alla camera \\
Questo perche d.c. è parito confessionale e ha appoggio di Chiesa e USA \\
Socialisti e comunisti si erano uniti nel fronte popolare, e insieme avevano fatto solo 31 per cento dei voti \\
Si apre nuova guerra tra comunismo e anticomunismo \t guerra civile, che ripropone a livello nazionale quello che accade a livello internazionale \t guerra fredda ha ripercussioni nei paesi \\
C'era spaccatura forte tra comunismo e anticomunismo \\
Il 14 luglio del 48 c'è attentato a Togliatti, che viene ferito da Pallante \t conservatore di destra, e da il via cosi all'occupazione della FIAt a torino, prendono in ostaggio dell'amministratore delegato e vogliono fare rivoluzione \\
Togliatti pero dal letto di ospedale dice fermatevi, fare insurrezione armata sarebbe tragico errore \t porterebbe l'italia in guerra civile \t togliatti richiama ad atteggiamento moderato \\
Tensioni hanno ripercussioni a livello sindacale \t la corrente cattolica esce dalla CGL, e fonda la CISL (confederazione italiana sindacati del lavoratori) \\
Viene fondata a Roma anche UIL (repubblicani e socialdemocratici, collocazione a meta strada tra CGL comunista e CISL d.c.) \\
Nel 1950 nasce anche la CISNAL \t chiaramente fascista, vicina all'MSI (movimento sociale italiano) = partito nato nel 46 che si richiama agli ideali del fascismo \\
Perche viene lasciato questo partito? \t perche democrazia, e partito si distanzia dalla prassi politica \t quindi viene ammesso

\v

Centrismo fa fino al 62 \t governa d.c. con il supporto di partiti di centro \t il centrismo va a sostituire i governi monocolore (dove governa un unico partito), ma in realta il partito forte era d.c., che dava impronta \\
D.c. governa con un programma di riforme e di mediazione \t il centrismo voleva essere un governo cautamente riformista (non coservatore), ma neanche troppo riformista \\
Le riforme pero non tengono il passo dell'evoluzione economico-sociale dell'italia \t a partire dal 58 c'è il "Miracolo economico" fino al 63, con boom economico \\
Italia in questo periodo diventa uno dei paesi piu industrializzati \t il via viene dato nel 58 quando entra nella CEE \t mentre nel 63 nasce il primo governo di centro sinistra \\
Il governo pero non sta al passo della crescita, non fanno riforme adeguate a questo cambiamento epocale x italia \\
Quei governi erano troppo vecchi, troppo conservatori \t e quindi il primo che apre alla possibilita a governi di centro-sinistra è Aldo Moro \\
La sua linea è quella di collaborare con i socialisti per il paese \t si apre quindi un epoca nuova da 62 a 68, che epoca del centro sinistra \t nascono governi misti con forze anche di sinistra, pero non comuniste \t solo socialdemocrazia, socialisti
\end{document}
